\chapter{Les contrats à terme et leurs bourses}\label{ch:m1:futures-contracts-and-exchanges}

Avec un indice à 6\,000, un \omterm{def:m1:futures-contracts-and-exchanges:future}{contrat à terme} E-mini S\&P 500 représente une exposition à
300\,000~\$ d'actions américaines. Il ne coûte rien à l'entrée~: personne ne paie
300\,000~\$, et le dépôt de garantie exigé par la chambre de compensation ne représente
qu'une petite fraction de ce montant. Le contrat peut être vendu à découvert aussi facilement qu'acheté, il se négocie du dimanche
soir au vendredi après-midi, et son plus petit pas de cotation, un quart de
point, vaut 12.50~\$. Chacun de ces faits figure dans un règlement de marché,
et il diffère pour chacun des milliers de \omterm{def:m1:futures-contracts-and-exchanges:future}{contrats à terme} cotés dans le monde.
Une société de trading qui négocie des \omterm{def:m1:futures-contracts-and-exchanges:future}{contrats à terme} commence par transformer les règlements de marché en un tableau. Ce
chapitre construit ce tableau~; les trois suivants expliquent comment les ordres s'apparient, comment le
dépôt de garantie est calculé, et comment le cours d'un \omterm{def:m1:futures-contracts-and-exchanges:future}{contrat à terme} se rattache
au sous-jacent sur lequel il porte.

\section{Le contrat}

\begin{definition}[Contrat à terme]\label{def:m1:futures-contracts-and-exchanges:future}
Un \emph{contrat à terme}\index{contrat à terme} est un accord standardisé,
coté par une bourse et garanti par sa chambre de compensation
(\cref{ch:m1:clearing-and-settlement}), pour acheter ou vendre une quantité déterminée
d'un sous-jacent à une date future déterminée, à un prix convenu aujourd'hui. Les gains et
pertes sont réglés en espèces chaque jour~; à l'échéance, le contrat est soit
réglé en espèces sur la base d'un cours de référence, soit livré par la remise du sous-jacent.
\end{definition}

La standardisation fait le produit. Deux parties qui veulent un \omterm{def:m1:futures-contracts-and-exchanges:future}{contrat à terme} de gré à gré sur le pétrole brut
doivent s'accorder sur la qualité, le lieu, la date, la taille, le crédit et la documentation. La
bourse fixe tous ces paramètres une fois pour toutes, de sorte que le seul élément restant à négocier est
le prix, et que tous les contrats d'un mois donné sont interchangeables
entre eux~: une position achetée auprès d'une contrepartie peut être clôturée en vendant à
une autre.

\begin{definition}[Multiplicateur, tick et valeur du tick]\label{def:m1:futures-contracts-and-exchanges:tick}
Le \emph{multiplicateur du contrat}\index{multiplicateur du contrat} convertit le cours
coté en argent~: le notionnel du contrat est égal au cours multiplié par le multiplicateur.
Le \emph{tick} est l'incrément minimal de cours, et la \emph{valeur du
tick}\index{valeur du tick} est égale au tick multiplié par le multiplicateur, soit le gain ou
la perte d'un contrat pour un tick.
\end{definition}

\begin{definition}[Code du mois d'échéance et échéance la plus proche]\label{def:m1:futures-contracts-and-exchanges:code}
Le symbole d'un \omterm{def:m1:futures-contracts-and-exchanges:future}{contrat à terme} est une racine de produit suivie d'un \emph{code du mois
d'échéance}\index{code du mois d'échéance}, une lettre pour le mois (F, G, H, J, K, M, N,
Q, U, V, X, Z pour janvier à décembre), puis le dernier chiffre de l'année. La
\emph{échéance la plus proche}\index{echeance@échéance la plus proche} est l'échéance cotée la plus rapprochée, généralement
la plus négociée.
\end{definition}

\begin{omfigure}
\begin{tikzpicture}[font=\small]
\node[font=\Huge\ttfamily,text=omDef] (r) at (0,0.2) {ES};
\node[font=\Huge\ttfamily,text=omThm] (m) at (1.7,0.2) {Z};
\node[font=\Huge\ttfamily,text=omExo] (y) at (3.1,0.2) {6};
\node[below=4pt of r,font=\footnotesize,align=center,text=omDef] {racine~:\\le produit};
\node[below=4pt of m,font=\footnotesize,align=center,text=omThm] {mois~:\\décembre};
\node[below=4pt of y,font=\footnotesize,align=center,text=omExo] {année~:\\2026};
\foreach \c/\n [count=\i] in {F/janv.,G/févr.,H/mars,J/avr.,K/mai,M/juin,N/juil.,Q/août,U/sept.,V/oct.,X/nov.,Z/déc.}{
  \pgfmathsetmacro{\x}{4.2+0.78*\i}
  \pgfmathtruncatemacro{\q}{mod(\i,3)==0}
  \ifnum\q=1
    \node[draw=omThm,fill=omThm!12,minimum width=0.7cm,minimum height=0.7cm,font=\ttfamily\bfseries] at (\x,0.25) {\c};
  \else
    \node[draw=omExo,minimum width=0.7cm,minimum height=0.7cm,font=\ttfamily] at (\x,0.25) {\c};
  \fi
  \node[font=\scriptsize] at (\x,-0.4) {\n};}
\end{tikzpicture}

\omcaption{Lecture d'un symbole de \omterm{def:m1:futures-contracts-and-exchanges:future}{contrat à terme}. Les codes grisés, H, M, U et Z, forment le
cycle trimestriel sur lequel sont cotés les \omterm{def:m1:futures-contracts-and-exchanges:future}{contrats à terme} sur indices actions et sur obligations~; les
contrats sur énergie sont cotés chaque mois.}\label{fig:m1:futures-contracts-and-exchanges:symbol}
\end{omfigure}

\section{Les groupes de bourses}

La liquidité des \omterm{def:m1:futures-contracts-and-exchanges:future}{contrats à terme} est concentrée par produit, et chaque produit est coté sur une seule
bourse~: contrairement aux actions (\cref{ch:m1:us-equity-market-structure}), un \omterm{def:m1:futures-contracts-and-exchanges:future}{contrat à terme}
appartient à la bourse qui le cote et ne se compense qu'auprès
de la chambre de compensation de cette bourse. Une position ouverte sur une bourse ne peut pas être
clôturée sur une autre. Il en résulte un petit nombre de groupes, chacun en
situation de quasi-monopole sur ses produits~: aux États-Unis, un groupe possède les
bourses de Chicago et de New York qui cotent les contrats sur indices actions de référence,
sur taux d'intérêt, énergie, métaux et matières agricoles~; un groupe atlantique
cote le pétrole Brent et les taux courts européens à Londres, ainsi que les matières premières molles
à New York~; le groupe boursier allemand cote les
\omterm{def:m1:futures-contracts-and-exchanges:future}{contrats à terme} sur indices actions et sur obligations d'État de la zone euro~; les bourses nationales d'Asie et
d'Amérique latine cotent leurs propres indices, taux et matières premières. La concurrence
entre bourses s'exerce au lancement d'un produit, et rarement par la suite.

\section{Six spécifications}

\begin{center}
\footnotesize
\begin{tabular}{llllll}
\toprule
Contrat & Unité & \omterm{def:m1:futures-contracts-and-exchanges:tick}{Tick} & \omterm{def:m1:futures-contracts-and-exchanges:tick}{Valeur du tick} & Échéances & Règlement \\
\midrule
E-mini S\&P 500 (ES) & 50~\$ $\times$ indice & 0.25 & 12.50~\$ & trimestrielles & espèces \\
EURO STOXX 50 (FESX) & 10~\euro $\times$ indice & 1 & 10~\euro & trimestrielles & espèces \\
T-note à 10 ans (ZN) & 100\,000~\$ nominal & $1/64$ & 15.625~\$ & trimestrielles & livraison \\
Euro-Bund (FGBL) & 100\,000~\euro nominal & 0.01 & 10~\euro & trimestrielles & livraison \\
WTI crude (CL) & 1\,000 barils & 0.01~\$ & 10~\$ & mensuelles & livraison \\
Brent crude (B) & 1\,000 barils & 0.01~\$ & 10~\$ & mensuelles & espèces ou EFP \\
\bottomrule
\end{tabular}
\end{center}

\begin{dated}{2026-09}{Détails tirés des règlements de marché}\label{dat:m1:futures-contracts-and-exchanges:specs}
\textbf{ES}~: le règlement-livraison final se fait en espèces sur la base d'un cours spécial d'ouverture (SOQ)
de l'indice, calculé à partir des cours d'ouverture de ses actions composantes le
troisième vendredi du mois d'échéance~; les calendar spreads se négocient par incréments
de 0.05. \textbf{FESX}~: les douze échéances trimestrielles les plus proches sont cotées~;
la négociation prend fin à 12:00 CET le troisième vendredi et le cours de règlement-livraison final
est la moyenne de l'indice entre 11:50 et 12:00~; la séance régulière va de
02:10 à 22:00 CET. \textbf{ZN}~: les notes livrables ont une maturité résiduelle
d'au moins six ans et demi et inférieure à huit ans~; les cours sont exprimés
en points et en trente-deuxièmes. \textbf{FGBL}~: une obligation notionnelle à 6~\%~;
les obligations livrables ont entre 8.5 et 10.5 ans de maturité résiduelle. \textbf{CL}~: livraison à
Cushing, Oklahoma. \textbf{B}~: jusqu'à 156 mois consécutifs cotés~; la négociation
prend fin le dernier jour ouvré du deuxième mois précédant le mois d'échéance~;
règlement-livraison par exchange for physical (EFP), avec une option de règlement en espèces sur
un indice.
\end{dated}

Les \omterm{def:m1:futures-contracts-and-exchanges:future}{contrats à terme} sur obligations se cotent dans une base qui précède l'informatique. Un cours de
\texttt{112-165} signifie $112 + 16.5/32$~: les deux premiers chiffres après le tiret
comptent des trente-deuxièmes de point, et un troisième chiffre, 2, 5 ou 7, un quart, une
moitié ou trois quarts de trente-deuxième. Avec un \omterm{def:m1:futures-contracts-and-exchanges:tick}{tick} d'une
demi-trente-deuxième, soit $1/64$ de point sur 100\,000~\$ de valeur nominale, la \omterm{def:m1:futures-contracts-and-exchanges:tick}{valeur du tick}
vaut 15.625~\$. Tout système qui manipule ces cours doit les représenter comme
des nombres entiers de \omterm{def:m1:futures-contracts-and-exchanges:tick}{ticks}~: un cours n'est alors jamais «~presque~» sur la grille, les tests d'égalité
sont exacts et aucun P\&L ne dépend d'un arrondi. (Un soixante-quatrième
est exactement représentable en virgule flottante binaire~; 0.01, le \omterm{def:m1:futures-contracts-and-exchanges:tick}{tick} de trois autres contrats
du tableau, ne l'est pas.)

\begin{omfigure}
\begin{tikzpicture}
\begin{axis}[width=10cm,height=4.6cm,ybar,bar width=14pt,
  ylabel={un tick (pb de valeur)},
  symbolic x coords={ES,FESX,ZN,FGBL,CL,B},xtick=data,
  tick label style={font=\small},label style={font=\small},xticklabel style={font=\small\ttfamily\strut},
  ymin=0,ymax=2.3,ymajorgrids,grid style={gray!25},
  nodes near coords,nodes near coords style={font=\footnotesize,/pgf/number format/.cd,fixed,fixed zerofill,precision=2},
  table/col sep=comma]
\addplot[fill=omDef!60,draw=omDef] table[x=root,y=tick_bp]{figdata/markets-1/18-futures-contracts-and-exchanges/specs.csv};
\end{axis}
\end{tikzpicture}

\omcaption{La taille du \omterm{def:m1:futures-contracts-and-exchanges:tick}{tick}, en fraction de la valeur du contrat, à des cours illustratifs
(ES 6\,000, FESX 5\,400, ZN 112, FGBL 128, CL 70, B 74). Un \omterm{def:m1:futures-contracts-and-exchanges:tick}{tick} large
implique une fourchette minimale élevée, de longues files d'attente et une prime à la priorité~; un
\omterm{def:m1:futures-contracts-and-exchanges:tick}{tick} fin favorise la concurrence par les prix. Les conséquences de cet appariement sont traitées
au \cref{ch:m1:matching-algorithms-and-implied-spreads}. Données~: le tableau du
tutoriel.}\label{fig:m1:futures-contracts-and-exchanges:tickbp}
\end{omfigure}

\section{Séances, règlement-livraison et limites}

\begin{definition}[Cours de compensation quotidien]\label{def:m1:futures-contracts-and-exchanges:settle}
Le \emph{cours de compensation quotidien}\index{cours de compensation quotidien} est le cours
publié par la bourse chaque jour, auquel chaque position ouverte est valorisée
et la \omterm{def:m1:clearing-and-settlement:margin}{marge de variation} calculée. Il est déterminé par une règle, typiquement une
moyenne pondérée par les volumes des transactions sur une courte fenêtre, et non par la dernière transaction.
\end{definition}

\begin{definition}[Positions ouvertes]\label{def:m1:futures-contracts-and-exchanges:oi}
Les \emph{positions ouvertes}\index{positions ouvertes} désignent le total de tous les contrats
conclus et non encore soldés par une transaction, par une livraison ou par
un exercice. Chaque contrat ouvert comporte un acheteur et un vendeur~; une transaction augmente
les positions ouvertes lorsque les deux côtés ouvrent, les réduit lorsque les deux côtés clôturent, et les laisse
inchangées lorsqu'un côté cède sa position à l'autre.
\end{definition}

\begin{definition}[Limite de variation de cours]\label{def:m1:futures-contracts-and-exchanges:limit}
Sur un \omterm{def:m1:futures-contracts-and-exchanges:future}{contrat à terme}, une \emph{\omterm{def:m1:asian-and-emerging-equity-markets:limit}{limite de variation de cours}} (pour les actions,
\cref{def:m1:asian-and-emerging-equity-markets:limit}) est une borne, fixée par la
bourse par rapport à un cours de référence, au-delà duquel le contrat ne peut
pas être négocié pendant une séance ou une partie de séance~; l'atteindre peut déclencher une suspension.
\end{definition}

\begin{dated}{2026-09}{Comment l'ES se règle chaque jour et quand il s'arrête}\label{dat:m1:futures-contracts-and-exchanges:limits}
Le règlement quotidien de l'E-mini de l'\omterm{def:m1:futures-contracts-and-exchanges:code}{échéance la plus proche} est le cours moyen pondéré par les volumes
des transactions entre 14:59:30 et 15:00:00 heure de Chicago, avec le
milieu entre le cours acheteur et le cours vendeur en repli~; les autres échéances sont réglées à partir des
cours de calendar spread. Des limites de baisse de 7~\%, 13~\% et 20~\% par rapport à un
cours de référence s'appliquent pendant les heures de marché des actions américaines, en coordination avec les
coupe-circuits du marché actions (\cref{ch:m1:exchange-traded-funds}), et une
bande symétrique s'applique la nuit (5~\% dans les deux sens dans le dépôt de 2016
consulté~; il faut consulter la règle en vigueur).
\end{dated}

Comme toutes les échéances d'un produit se règlent à quinze heures et qu'une position
y est valorisée, les trente dernières secondes avant le règlement sont, pour un
desk de \omterm{def:m1:futures-contracts-and-exchanges:future}{contrats à terme}, ce que la criée de clôture est pour un desk actions.

Le volume passe d'une échéance à la suivante dans les quelques jours qui précèdent l'échéance~: le
\emph{roll}. Les détenteurs qui veulent conserver leur exposition vendent l'\omterm{def:m1:futures-contracts-and-exchanges:code}{échéance la plus proche} et achètent la
suivante sous forme de calendar spread~; les \omterm{def:m1:futures-contracts-and-exchanges:oi}{positions ouvertes} migrent, et l'\omterm{def:m1:futures-contracts-and-exchanges:code}{échéance la plus
proche} des données change. Une série de cours raccordée à travers les roll sans ajustement
contient des sauts qui ne sont pas des rendements
(\cref{ch:m1:basis-roll-and-delivery}).

\begin{omfigure}
\begin{tikzpicture}
\begin{axis}[width=11cm,height=4.8cm,
  xlabel={jours de bourse avant l'échéance du contrat le plus proche},ylabel={positions ouvertes (millions)},
  tick label style={font=\small},label style={font=\small},
  xmin=-30,xmax=9,ymin=0,ymax=2.8,grid=major,grid style={gray!25},
  legend columns=2,legend style={font=\footnotesize,at={(0.5,-0.32)},anchor=north,draw=none,fill=none,/tikz/every even column/.append style={column sep=8pt}},
  table/col sep=comma]
\addplot[omDef,very thick] table[x=day,y=front_m]{figdata/markets-1/18-futures-contracts-and-exchanges/roll.csv};
\addplot[omThm,very thick,dashed] table[x=day,y=next_m]{figdata/markets-1/18-futures-contracts-and-exchanges/roll.csv};
\legend{échéance la plus proche,échéance suivante}
\end{axis}
\end{tikzpicture}

\omcaption{Un roll trimestriel stylisé~: les \omterm{def:m1:futures-contracts-and-exchanges:oi}{positions ouvertes} passent de l'\omterm{def:m1:futures-contracts-and-exchanges:code}{échéance la plus proche}
à la suivante en une semaine environ, centré ici à huit jours de bourse avant l'échéance. La date et
la vitesse diffèrent selon le produit et doivent être mesurées, et non supposées. Données~: la
simulation du tutoriel.}\label{fig:m1:futures-contracts-and-exchanges:roll}
\end{omfigure}

\section{Qui négocie les contrats à terme}

\begin{definition}[Hedger et spéculateur]\label{def:m1:futures-contracts-and-exchanges:hedger}
Un \emph{hedger}\index{hedger} détient des \omterm{def:m1:futures-contracts-and-exchanges:future}{contrats à terme} pour compenser un risque qu'il porte
par ailleurs~: la récolte d'un producteur, les actions d'un fonds, les obligations d'un dealer. Un
\emph{spéculateur}\index{spéculateur} les détient pour l'exposition elle-même. Le rapport hebdomadaire
du régulateur américain sur les positions classe les grands intervenants en
\emph{commerciaux}, ceux qui utilisent le contrat pour se couvrir selon la
définition de la réglementation, et en \emph{non-commerciaux}.
\end{definition}

Le couple théorique masque la population qui compte pour une société de trading~:
les \omterm{def:m1:what-a-trading-firm-does:asset-manager}{sociétés de gestion} qui convertissent du cash en actions et gèrent la duration~; les fonds à effet de levier dans
les opérations de base, de tendance et de valeur relative~; les dealers qui couvrent les swaps et les options~;
et les \omterm{def:m1:what-a-trading-firm-does:market-maker}{teneurs de marché}, qui conservent peu de positions la nuit et représentent une grande part des
volumes. Un \omterm{def:m1:futures-contracts-and-exchanges:future}{contrat à terme} est aussi la couverture de choix de tout instrument \omterm{def:m1:delta-one-instruments:deltaone}{delta-one}
du \cref{ch:m1:delta-one-instruments}, qui lie son cours au marché au comptant
dans les bornes du \cref{ch:m1:basis-roll-and-delivery}.

\begin{method}[Dimensionnement d'une couverture par contrat à terme]\label{met:m1:futures-contracts-and-exchanges:hedge}
Pour couvrir un portefeuille de valeur $V$ et de bêta $\beta$ par rapport à l'indice au moyen d'un \omterm{def:m1:futures-contracts-and-exchanges:future}{contrat à terme}
de cours $F$ et de multiplicateur $m$~: vendre $N = \beta V/(mF)$ contrats, arrondi
à l'entier le plus proche. L'exposition résiduelle est $\beta V - NmF$, au plus
égale à une demi-taille de contrat~; pour les petits portefeuilles, utiliser le contrat le plus petit s'il existe.
Le bêta est une estimation~: son erreur est généralement bien supérieure à l'arrondi.
\end{method}

\begin{omfigure}
\begin{tikzpicture}
\begin{axis}[width=11cm,height=4.8cm,
  xlabel={valeur du portefeuille (millions de \$)},ylabel={résiduel non couvert (\%)},
  tick label style={font=\small},label style={font=\small},
  xmin=0.1,xmax=3,ymin=0,ymax=55,grid=major,grid style={gray!25},
  legend columns=2,legend style={font=\footnotesize,at={(0.5,-0.32)},anchor=north,draw=none,fill=none,/tikz/every even column/.append style={column sep=8pt}},
  table/col sep=comma]
\addplot[omDef,thick] table[x=portfolio_m,y=residual_pct_es]{figdata/markets-1/18-futures-contracts-and-exchanges/residual.csv};
\addplot[omThm,thick,dashed] table[x=portfolio_m,y=residual_pct_micro]{figdata/markets-1/18-futures-contracts-and-exchanges/residual.csv};
\legend{contrat de 300\,000~\$,contrat dix fois plus petit}
\end{axis}
\end{tikzpicture}

\omcaption{Arrondi aux contrats entiers~: le résiduel d'exposition d'un
portefeuille de bêta un, en pourcentage de sa valeur. En dessous de 150\,000~\$, le grand
contrat ne peut pas couvrir du tout (résiduel de 100~\%, hors de l'échelle). Données~:
\cref{met:m1:futures-contracts-and-exchanges:hedge}.}\label{fig:m1:futures-contracts-and-exchanges:residual}
\end{omfigure}

\section{Tutoriel~: un tableau de contrats}\label{tut:m1:futures-contracts-and-exchanges:table}

\begin{tutorial}
\textbf{Objectif.} Conserver six spécifications de contrats dans une même structure et en déduire
le notionnel, la \omterm{def:m1:futures-contracts-and-exchanges:tick}{valeur du tick}, le \omterm{def:m1:futures-contracts-and-exchanges:tick}{tick} en points de base, l'effet de levier et les tailles de
couverture à partir de celle-ci. \textbf{État final~:} les trois figures de données de ce chapitre.
\begin{tutsteps}
\item \textbf{Un enregistrement par contrat.} Les multiplicateurs des obligations s'expriment par
point de cotation~: 100\,000~\$ de valeur nominale cotés en pourcentage valent 1\,000~\$ le point.
\omcode{code/markets-1/18-futures-contracts-and-exchanges/python/futures_basics.py}{8}{27}{Une spécification et les trois nombres qui en découlent.}\label{lst:m1:futures-contracts-and-exchanges:spec}
\item \textbf{Dimensionnement de la couverture.} Arrondir au contrat le plus proche et indiquer
ce qui reste.
\omcode{code/markets-1/18-futures-contracts-and-exchanges/python/futures_basics.py}{44}{48}{Nombre de contrats à vendre et exposition résiduelle.}\label{lst:m1:futures-contracts-and-exchanges:hedge}
\item \textbf{Effet de levier.} Diviser le notionnel par la \omterm{def:m1:clearing-and-settlement:margin}{marge initiale} que propose le courtier
aujourd'hui (\cref{ch:m1:margin})~; avec une marge de 5~\% du notionnel, un mouvement de 2~\%
représente 40~\% du dépôt.
\end{tutsteps}
\textbf{Pour aller plus loin.} Ajouter quatre contrats tirés des pages produits des bourses,
parmi eux un \omterm{def:m1:futures-contracts-and-exchanges:future}{contrat à terme} sur taux court, dont le cours vaut
100 moins un taux et dont la \omterm{def:m1:futures-contracts-and-exchanges:tick}{valeur du tick} provient d'une convention de décompte, non
d'un cours.
\end{tutorial}

\section{Construction~: le référentiel des contrats}\label{bld:m1:futures-contracts-and-exchanges:master}

\begin{build}
\bfield{Rôle} Chaque composant de la petite société fictive qui touche
à un \omterm{def:m1:futures-contracts-and-exchanges:future}{contrat à terme} (le tenancier des P\&L du \cref{ch:m1:pnl-and-positions}, le moteur de marges
du \cref{ch:m1:margin}, la saisie des ordres) demande à un seul module ce qu'est le contrat.
\bfield{Interface} \texttt{ContractMaster.from\_csv(path)};
\texttt{spec(root)} avec \texttt{multiplier}, \texttt{tick\_size},
\texttt{tick\_value} exacts, \texttt{to\_ticks(price)}, \texttt{from\_ticks(n)},
\texttt{notional(price)}~; \texttt{parse(symbol, today)} renvoyant un
instrument~; \texttt{expiry(instrument)}~;
\texttt{parse\_thirty\_seconds(text)}.
\bfield{Règles} Tous les calculs en fractions exactes~; un cours hors de la grille de \omterm{def:m1:futures-contracts-and-exchanges:tick}{tick} est une
erreur, jamais arrondi. Une année à un seul chiffre se résout en la première de ces années
qui n'est pas passée. Un code de mois non prévu pour le produit est une
erreur. Les règles d'échéance qui exigent un calendrier de bourse lèvent une exception plutôt que de deviner.
\bfield{Données} \texttt{data/markets-1/contracts\_sample.csv}, six lignes
recopiées depuis les sources de ce chapitre~: un jeu de test, pas des données de référence.
\bfield{Tests d'acceptation} \texttt{code/firm/contracts/tests/}~: valeurs exactes des
\omterm{def:m1:futures-contracts-and-exchanges:tick}{ticks}~; allers-retours par des \omterm{def:m1:futures-contracts-and-exchanges:tick}{ticks} entiers~; cotations de Treasury~; symboles et
troisièmes vendredis~; validation.
\bfield{Pour aller plus loin} Charger le fichier quotidien de données de référence d'une bourse et
le comparer à celui de la veille~: les changements de pas de cotation et d'admission à la cote sont des événements
(\cref{ch:m1:reading-market-data}).
\end{build}

\begin{omsources}
\item CME, Rulebook Chapter 358 (E-mini S\&P 500 futures) et Chapter 351,
tel que déposé auprès de la CFTC (mars 2016)~; CME Group, daily settlement procedures
for equity-index futures, tel que déposé auprès de la CFTC (janvier 2018).
\item CBOT, Rulebook Chapter 19 (Treasury note futures), et le dépôt auprès de la CFTC
du 21 février 2025 sur les qualités livrables~; NYMEX, Rulebook Chapter 200.
\item Eurex, pages produits \emph{EURO STOXX 50 Index Futures} et
\emph{Euro-Bund Futures}~; ICE Futures Europe, \emph{Brent Crude Futures}
page produit.
\item US Commodity Futures Trading Commission, \emph{Commitments of Traders:
Explanatory Notes}.
\end{omsources}

\section{Exercices}

\begin{exercise}[$\star$]\label{exo:m1:futures-contracts-and-exchanges:1}
Un trader est acheteur de 12 E-mini depuis 6\,012.50~; le cours baisse à 5\,987.25. Donner
le mouvement en \omterm{def:m1:futures-contracts-and-exchanges:tick}{ticks} et la perte.
\end{exercise}

\begin{exercise}[$\star$]\label{exo:m1:futures-contracts-and-exchanges:2}
Le \omterm{def:m1:futures-contracts-and-exchanges:future}{contrat à terme} sur obligation à 10 ans passe de \texttt{111-245} à \texttt{112-080}. Donner
les deux cours en décimal, le mouvement en \omterm{def:m1:futures-contracts-and-exchanges:tick}{ticks} et le profit sur un
contrat acheté.
\end{exercise}

\begin{exercise}[$\star$]\label{exo:m1:futures-contracts-and-exchanges:3}
Nous sommes le 18 septembre 2026. Décoder \texttt{FESXH7} et donner son dernier jour de
négociation. Que signifierait \texttt{ESH6} aujourd'hui~?
\end{exercise}

\begin{exercise}[$\star\star$]\label{exo:m1:futures-contracts-and-exchanges:4}
D'après la \cref{fig:m1:futures-contracts-and-exchanges:tickbp}, le \omterm{def:m1:futures-contracts-and-exchanges:tick}{tick} du FESX est
plus de quatre fois celui de l'ES, rapporté à la valeur. Donner deux conséquences
pour un \omterm{def:m1:what-a-trading-firm-does:market-maker}{teneur de marché} et une pour un fonds qui traverse la fourchette sur 50~millions d'euros.
\end{exercise}

\begin{exercise}[$\star\star$]\label{exo:m1:futures-contracts-and-exchanges:5}
Un courtier demande 21\,000~\$ de \omterm{def:m1:clearing-and-settlement:margin}{marge initiale} par E-mini à 6\,000. Donner
l'effet de levier et le mouvement de l'indice qui fait perdre tout le dépôt.
Pourquoi le dépôt n'est-il pas la bonne mesure de l'argent nécessaire pour tenir la position~?
\end{exercise}

\begin{exercise}[$\star\star$]\label{exo:m1:futures-contracts-and-exchanges:6}
Un nouveau contrat est admis à la cote. Transactions~: A achète 5 à B~; C achète 3 à A~; B achète 2
à C~; D achète 4 à A. Donner les \omterm{def:m1:futures-contracts-and-exchanges:oi}{positions ouvertes} et la position de chaque trader
après chaque transaction. Quels ont été le volume et les \omterm{def:m1:futures-contracts-and-exchanges:oi}{positions ouvertes} en fin de séance~?
\end{exercise}

\begin{exercise}[$\star\star\star$]\label{exo:m1:futures-contracts-and-exchanges:7}
\emph{Calcul.} À partir de \texttt{residual.csv}~: pour les portefeuilles de 1~million de dollars
et plus, indiquer le plus grand résidu avec le contrat de 300\,000~\$ puis avec le contrat
dix fois plus petit. Énoncer la borne générale et la vérifier.
\end{exercise}

\begin{exercise}[$\star\star\star$]\label{exo:m1:futures-contracts-and-exchanges:8}
\emph{Trouver la faille.} Un commentaire affirme~: «~Les \omterm{def:m1:futures-contracts-and-exchanges:oi}{positions ouvertes} sur l'\omterm{def:m1:futures-contracts-and-exchanges:code}{échéance
la plus proche} ont augmenté de 20~\% aujourd'hui alors que le cours a monté~: de nouveaux
acheteurs affluent, et les acheteurs sont désormais plus nombreux que les vendeurs.~» Corriger
cette affirmation et dire ce que les mêmes données peuvent légitimement apprendre.
\end{exercise}

\section{Problème~: Dimensionner une couverture}

\begin{problem}[{Problème du week-end --- un portefeuille transatlantique et deux contrats à terme}]\label{pb:m1:futures-contracts-and-exchanges:1}
Le 18 septembre 2026, un fonds veut supprimer l'exposition au marché de deux
poches pendant trois mois~: 187~millions de dollars d'actions américaines de bêta 1.15 par rapport
au S\&P 500, et 43~millions d'euros d'actions de la zone euro de bêta 0.90 par rapport à
l'EURO STOXX 50. L'E-mini est à 6\,000 et le FESX à 5\,400.

\medskip\noindent\textbf{Partie I --- Les contrats.}
\begin{enumerate}
\item Donner le notionnel d'un contrat de chaque type.
\item Donner le nombre exact et le nombre arrondi d'E-mini à vendre, ainsi que
le résidu.
\item Même question pour le FESX.
\item Quelle échéance le fonds doit-il vendre, avec son symbole et son dernier jour
de négociation, s'il ne veut pas faire de roll~?
\item Avec une \omterm{def:m1:clearing-and-settlement:margin}{marge initiale} de 5~\% du notionnel, donner le dépôt pour la couverture
américaine.
\end{enumerate}

\medskip\noindent\textbf{Partie II --- Une mauvaise journée.}
Le S\&P 500 recule de 2~\%.
\begin{enumerate}[resume]
\item Donner le gain sur les \omterm{def:m1:futures-contracts-and-exchanges:future}{contrats à terme} et la perte attendue sur la poche américaine.
\item Quand les liquidités issues des \omterm{def:m1:futures-contracts-and-exchanges:future}{contrats à terme} arrivent-elles, et quand la perte sur
les actions est-elle réalisée~? Pourquoi la différence compte-t-elle lors d'une \emph{bonne} journée~?
\item La poche perd en réalité 2.9~\%. Donner le P\&L couvert et en nommer
la source.
\item L'estimation du bêta du fonds a un écart-type de 0.10. Comparer
l'exposition correspondante avec le résidu d'arrondi.
\end{enumerate}

\medskip\noindent\textbf{Partie III --- Détails.}
\begin{enumerate}[resume]
\item En décembre, le fonds prolonge la couverture d'un trimestre et fait un roll
vers mars. Le calendar spread a un \omterm{def:m1:futures-contracts-and-exchanges:tick}{tick} de 0.05. Donner le coût du paiement d'un \omterm{def:m1:futures-contracts-and-exchanges:tick}{tick}
de spread sur l'ensemble de la position.
\item La couverture européenne est en euros et le fonds publie en dollars. Quelle
exposition la couverture n'a-t-elle \emph{pas} supprimée~?
\item Le \omterm{def:m1:futures-contracts-and-exchanges:settle}{cours de compensation quotidien} de l'E-mini est une moyenne sur trente secondes à
15~h~00 heure de Chicago, soit 16~h~00 à New York~; les actions sont valorisées à leur enchère
de clôture au même moment. Qu'apparaît dans le P\&L quotidien~?
\item Pourquoi le fonds couvrirait-il la poche américaine avec de l'ES alors que la majeure
partie de celle-ci est constituée de petites capitalisations, et à quoi il renonce~?
\item La couverture sera en place pendant trois mois. En s'appuyant sur les résultats
du \cref{ch:m1:delta-one-instruments}, que coûte-t-elle ou que rapporte-t-elle en dehors
des commissions~?
\end{enumerate}

\medskip\noindent\textbf{Partie IV --- Jugement.}
\begin{enumerate}[resume]
\item Après la couverture, quels risques le fonds continue-t-il de porter~?
\item Une baisse de 7~\% déclenche la première \omterm{def:m1:asian-and-emerging-equity-markets:limit}{limite de variation de cours} pendant les heures
américaines. Que devient la couverture, et qu'en est-il des actions~?
\item Un contrat dix fois plus petit améliorerait-il la couverture ici~?
\item Pourquoi les fonds couvrent-ils avec des \omterm{def:m1:futures-contracts-and-exchanges:future}{contrats à terme} plutôt qu'en vendant les actions~?
\item Énoncer le \emph{résultat nommé}~: le nombre d'E-mini et
l'exposition résiduelle.
\item En une phrase~: que standardise un \omterm{def:m1:futures-contracts-and-exchanges:future}{contrat à terme}, et que
laisse-t-il au trader~?
\end{enumerate}
\end{problem}

\section{Questions d'entretien}

\begin{interviewq}[$\star$ \iqroles{trader, développeur, chercheur}]\label{iq:m1:futures-contracts-and-exchanges:1}
Quelle est la \omterm{def:m1:futures-contracts-and-exchanges:tick}{valeur du tick} d'un \omterm{def:m1:futures-contracts-and-exchanges:future}{contrat à terme} E-mini S\&P 500, et que
vaut le contrat~?
\end{interviewq}

\begin{interviewq}[$\star$ \iqroles{trader, chercheur}]\label{iq:m1:futures-contracts-and-exchanges:2}
Qu'est-ce que les \omterm{def:m1:futures-contracts-and-exchanges:oi}{positions ouvertes}, et en quoi diffèrent-elles du volume~?
\end{interviewq}

\begin{interviewq}[$\star\star$ \iqroles{développeur, chercheur}]\label{iq:m1:futures-contracts-and-exchanges:3}
Comment représenteriez-vous les cours des \omterm{def:m1:futures-contracts-and-exchanges:future}{contrats à terme} dans un système de trading, et pourquoi~?
\end{interviewq}

\begin{interviewq}[$\star\star$ \iqroles{trader, chercheur}]\label{iq:m1:futures-contracts-and-exchanges:4}
Pourquoi un \omterm{def:m1:futures-contracts-and-exchanges:future}{contrat à terme} n'est-il négocié que sur une seule bourse alors qu'une action se
négocie sur de nombreuses plateformes~?
\end{interviewq}

\begin{interviewq}[$\star\star$ \iqroles{chercheur, ingénieur ML}]\label{iq:m1:futures-contracts-and-exchanges:5}
On vous remet dix ans de cours quotidiens de l'\omterm{def:m1:futures-contracts-and-exchanges:code}{échéance la plus proche}. Que devez-vous
faire avant de calculer les rendements~?
\end{interviewq}

\begin{interviewq}[$\star\star\star$ \iqroles{trader, chercheur}]\label{iq:m1:futures-contracts-and-exchanges:6}
Vous devez couvrir 200~millions de dollars d'actions pour un trimestre. Détaillez-moi chaque
décision.
\end{interviewq}
